\cotefile{79}
\scanpage{81}
les structures $(\widetilde F_{M,i}\ \unsure{\text{de}\mathbin{+}\pi|\text{déf}\ \pi|)_{i\in\pi}|)_{i\in\Pi}}$.

Soit $u$ un autom. de cette structure,
commutant aux \deleted{\unsure{K}} $\widetilde T_\lambda$, et fixant un
repère $(x,y,z)$, i.e. fixant les éléments
d'une base $(x,y)$.
% slanted
\margin{il suffit pour voir qu'il fixe $x,y$} \deleted{Alors} Notons que un
tel élément $\xi$, \inserted{$\in\widetilde S_M=M^{**}$} il fixe les $T_\lambda\xi=\lambda\xi$, $\lambda\in\Lambda^*$,
en particulier aussi $-\xi$ -- \deleted{on voit donc}
\inserted{Comme il} \deleted{par le lemme p.~$37'$, que \unsure{pour} $a,b\in$}
\inserted{si $x$} fixe $x$ et les $\lambda y$ $(\lambda\in\Lambda^*)$, on voit que il
\deleted{induit l'identit\unsure{é}} fixe les $\tau_A(\{x,\lambda y\})=\lambda\cdot i$
$(\text{où}\ i=\deleted{\unsure{\partial}}\tau_A\{x,y\})$, donc
il induit l'identité sur $\pi_M$, i.e. \deleted{\unsure{Raisonnons}} \inserted{c'est un}
\deleted{\unsure{P}}autom. pour toutes les structures $\widetilde F_{M,i}$.

\deleted{En vertu du lemme p.~$37'$, avec $x,y$,}
\deleted{Il fixe aussi} Considérons les matrices
\deleted{$u$ fixe aussi les $ax+by$ avec $a$}
\[
\left\{\begin{aligned}
\underline{\lambda}_0(\lambda)&=\begin{pmatrix}1&-\lambda\\0&1\end{pmatrix}
&&\lambda\in\Lambda^*\\
\lambda_1&=\begin{pmatrix}0&1\\-1&0\end{pmatrix}\\
\lambda_\infty&=\begin{pmatrix}1&-1\\1&0\end{pmatrix}
\end{aligned}\right.
\]

Alors l'ens. des bases $(x,y)$ ``fixes sous $u$'' est stable par composition avec les matrices
précédentes et leurs inverses, donc aussi
les composés de telles. Si $\Lambda^*$ engendre $\Lambda$
additivement (p.~ex. $\Lambda$ $\frac12$ local) alors \deleted{\unsure{du}} être
% archive sheet number: 80
